% File: appendices/appK-complexity.tex

\chapter{Computational Complexity Background}
\label{app:complexity}

This appendix summarizes the classical theory of computational complexity
required for the formal development of SpherePop, semantic computation, and
distributed lamphrodynamic processes in Volume~III.  The treatment emphasizes
structural complexity (reductions, completeness, hierarchies) and those aspects
which reappear in homotopy--theoretic or geometric encodings.  We work
predominantly in the classical Turing model and remark on geometric
representations when appropriate.

\section{Deterministic and Nondeterministic Computation}

\begin{definition}[Turing machine]
A (deterministic) Turing machine $M$ consists of a finite set of states,
a finite tape alphabet with a distinguished blank symbol,
a transition function, and a distinguished halting configuration.
Given an input $x \in \{0,1\}^\ast$, $M$ either halts with an output or
continues indefinitely.
\end{definition}

\begin{definition}[Time complexity]
The (worst--case) time complexity of $M$ is the function
\[
T_M(n) := \max \{ \text{number of steps $M$ uses on input $x$} \;:\; |x|=n \}.
\]
We say that $M$ runs in time $T(n)$ if $T_M(n)\le T(n)$ for all~$n$.
\end{definition}

\begin{definition}[Complexity classes]
\begin{enumerate}
  \item ${\sf P}$ consists of all problems decidable in polynomial time
        by a deterministic Turing machine.
  \item ${\sf NP}$ consists of all problems decidable in polynomial time
        by a nondeterministic Turing machine.
  \item ${\sf PSPACE}$ consists of problems decidable using polynomial space.
\end{enumerate}
\end{definition}

\begin{remark}
The containment ${\sf P} \subseteq {\sf NP} \subseteq {\sf PSPACE}$ is known,
but no strict separation of these classes has been proved.
The open question ${\sf P} \stackrel{?}{=} {\sf NP}$ is central to classical
complexity theory.
\end{remark}

\section{Reductions and Completeness}

\begin{definition}[Polynomial--time reduction]
A language $A$ reduces to $B$ (written $A \le_p B$) if there exists a
polynomial--time computable function $f$ such that
$x \in A$ iff $f(x)\in B$.  A problem is {\sf NP--hard} if every problem in
${\sf NP}$ reduces to it, and {\sf NP--complete} if it is both
{\sf NP--hard} and in ${\sf NP}$.
\end{definition}

\begin{remark}
For geometric computation models, one typically constructs an encoding
$\mathrm{Enc}(x)$ of classical input $x$ as a geometric configuration such that
logical operations correspond to merge or collapse operations.
Polynomial simulation requires that the size and depth of the encoding scale
polynomially with~$|x|$.
\end{remark}

\section{Alternation and Space Complexity}

\begin{definition}[Alternating Turing machine]
An alternating Turing machine is a nondeterministic Turing machine whose
states are partitioned into \emph{existential} and \emph{universal} states,
with acceptance defined inductively on the computation tree.
\end{definition}

\begin{theorem}[Chandra--Kozen--Stockmeyer]
For $k\ge 1$,
\[
{\sf AP} = {\sf PSPACE}.
\]
More precisely, polynomial--time alternating computation equals polynomial
space deterministic computation.
\end{theorem}

\begin{remark}
This equivalence plays a role in geometric encodings where
``branching'' is represented topologically (e.g.~via choice of merge
sequencing), and universal quantification appears as a geometric constraint.
\end{remark}

\section{Space--Time Hierarchies}

\begin{theorem}[Time hierarchy]
For reasonable $t(n)$,
\[
\mathrm{DTIME}(t(n)) \subsetneq \mathrm{DTIME}(t(n)\log t(n)).
\]
\end{theorem}

\begin{theorem}[Space hierarchy]
For space--constructible $s(n)$,
\[
\mathrm{DSPACE}(s(n)) \subsetneq \mathrm{DSPACE}(s(n)\log s(n)).
\]
\end{theorem}

\begin{remark}
Such strict inclusions demonstrate that increasing resources truly increases
computational power in the Turing model; analogous results for geometric
computations require careful resource measures (size of configuration,
dimensionality, number of merge operations).
\end{remark}

\section{Circuit Complexity and Uniformity}

\begin{definition}[Boolean circuit]
A Boolean circuit is a directed acyclic graph of logical gates computing a
Boolean function $f:\{0,1\}^n\to \{0,1\}$.
Its size is the number of gates and its depth is the length of the longest
directed path.
\end{definition}

\begin{definition}[Circuit classes]
\begin{enumerate}
  \item ${\sf AC}^0$ are constant--depth, polynomial--size circuits with
        unbounded fan--in AND/OR gates.
  \item ${\sf NC}^k$ are poly--size circuits of depth $O(\log^k n)$.
\end{enumerate}
\end{definition}

\begin{remark}
Circuit classes give an alternative viewpoint for merge--collapse systems:
merge corresponds to parallel composition, while collapse relates to depth
reduction or abstraction.  Establishing lower bounds in this setting is an
open problem (see \cref{sec:open-problems-volume3}).
\end{remark}

\section{Encoding Lambda Calculus and Turing Machines}

\begin{definition}[Lambda calculus]
The untyped lambda calculus consists of variables, abstractions $\lambda x.M$,
and applications $M\,N$.  Reduction rules include $\beta$--reduction
$(\lambda x.M)\,N \to M[x:=N]$.
\end{definition}

\begin{theorem}[Church--Turing thesis (informal)]
The lambda calculus, Turing machines, and other standard models compute the
same class of (partial) functions on natural numbers.
\end{theorem}

\begin{remark}
In SpherePop, lambda terms will be represented as geometric configurations,
where application is modeled by merge, and abstraction corresponds to a
collapse or quotient identifying a boundary component.  Complexity bounds
require that the encoding and its evolution remain polynomially controlled.
\end{remark}

\section{Complexity in Geometric and Homotopy--Theoretic Models}

\begin{definition}[Size of a geometric configuration]
Fix a class of finite geometric objects (e.g.~triangulated manifolds, or
finite unions of balls).  The \emph{size} of a configuration is the number of
primitive pieces needed to describe it (cells, simplices, generators).
\end{definition}

\begin{remark}
A merge--collapse system is computationally \emph{efficient} if each elementary
operation transforms a configuration of size $n$ into one of size $O(n^c)$ for
some fixed~$c$, and the number of operations needed to simulate a Turing step
is polynomial in $n$.  This is a key requirement in the proof of Turing
completeness in \cref{chap:spherepop-computation}.
\end{remark}

\section{Distributed Complexity and CRDTs}

\begin{definition}[Distributed complexity (informal)]
Given a network of processes updating a shared state, the distributed
complexity measures the number of rounds or messages required to achieve
consistency or agreement under specified failure models.
\end{definition}

\begin{remark}
CRDTs (Conflict--Free Replicated Data Types) provide eventual consistency using
commutative operations.  Interpreting these operations as lamphrodynamic flows
introduces a geometric semantics in which convergence properties become
statements about dissipative field evolution.
\end{remark}

\section{Equivalence Problems and Higher Categories}

\begin{remark}
Determining equivalence of objects in an $\infty$--category is in general far
more complex than equality of classical data.  Complexity grows rapidly with
dimension and coherence conditions.  At present, only partial complexity
classifications are known; these motivate the conjectured hardness results in
Volume~III for geometric computation of equivalence.
\end{remark}

\section{Summary}

This appendix has outlined the classical foundations of computational
complexity relevant to geometric and homotopy--theoretic computation.  The key
concepts used in later chapters are reductions, completeness, complexity
classes, hierarchy theorems, and the geometric interpretation of
merge--collapse operations.  Further references are provided in the unified
bibliography.


